
\section{The Rotation Group}\label{sec:rotations}

\hypertarget{definition}{%
\subsection{Definition}\label{definition}}

Geometrically, a rotation is a transformation of a plane or space that causes no
deformation or reflection so that distance, orientation, and the origin
are preserved. Mathematically this is realized in \(\mathbb{R}^3\) by
considering \emph{all} functions
\(\bm{\Lambda}\colon \mathbb{R}^3\rightarrow \mathbb{R}^3\), and
filtering out those which do not preserve the origin \(\bm{0}\),
the vector dot product \(\bm{a}\cdot \bm{b}\),
and the scalar triple product
\([\bm{a},\bm{b},\bm{c}] = \bm{a}\cdot\left(\bm{b}\times\bm{c}\right)\)
for vectors \(\bm{a},\bm{b}\) and \(\bm{c}\) in \(\mathbb{R}^3\). The
set \(\mathrm{SO}(3)\) is defined as the collection of all functions
that have these properties: %features:
\begin{equation}
\mathrm{SO}(3) \triangleq \left\{\,
    \bm{\Lambda}\colon \mathbb{R}^3\rightarrow\mathbb{R}^3
~\left|~~
\begin{aligned}
\bm{\Lambda}\bm{a}\cdot\bm{\Lambda}\bm{b} &= \bm{a}\cdot\bm{b}\,
,  \\
[\bm{\Lambda a}, \bm{\Lambda b}, \bm{\Lambda c}]&=[\bm{a}, \bm{b}, \bm{c}]
,~\text{ and } ~\bm{\Lambda}\bm{0} = \bm{0}
\text{ for all }\bm{a},\bm{b},\bm{c}\in\mathbb{R}^3
\end{aligned}
\right.
\,\right\}
\label{eq:def-SO3}\end{equation}
These properties alone that ensure any such map \(\bm{\Lambda}\) acts linearly on vectors\footnote{The linearity of $\mathrm{SO}(n)$ is more commonly specified directly by taking it as a subgroup of the general linear group, $\mathrm{GL}(n)$ (see e.g. \cite{artin1991algebra}).
However, it is helpful to observe this implication from the rotation constraints in Equation~(\ref{eq:def-SO3}).}. 
The properties in Equation~(\ref{eq:def-SO3}) for any 
\(\bm{a},\bm{b},\bm{c}\in\mathbb{R}^3\)
imply the familiar relations:
\begin{equation*} 
\begin{aligned}
\bm{\Lambda}\bm{a}\cdot\bm{\Lambda}\bm{b} &= \bm{a}\cdot\bm{b}\,
  \\
[\bm{\Lambda a}, \bm{\Lambda b}, \bm{\Lambda c}]&=[\bm{a}, \bm{b}, \bm{c}] \\
\end{aligned} \ \Longleftrightarrow
\begin{aligned}
\bm{\Lambda}^{\mathrm{t}} \bm{\Lambda}&=\mathbf{1} \\
\det\bm{\Lambda} &= 1 .
\end{aligned}
\end{equation*}
%where the left-hand side is asserted again for any
%\(\bm{a},\bm{b},\bm{c}\in\mathbb{R}^3\).

\hypertarget{representation}{%
\subsection{Representation}\label{representation}}

In this work, the term \emph{representation} is used to indicate how a
rotation \(\bm{\varLambda}\) is stored in computer memory, and how
its action on vectors and matrices is computed.
This is unrelated to the \emph{parameterization} of a given problem and is not considered a transformation of the configuration space.
Any of the standard ``parameterizations'' can be used both to \emph{represent}
rotation operations, and to \emph{parameterize} a configuration space.
Well known examples %\FCF{alternatives} %examples
% NOTE: "examples" is used in the sense of "examples of the standard parameterizations mentioned in the preceeding sentence". I dont see what these are "alternatives" to.
include Euler angles, and the Euler, Rodriguez, and Cayley parameters.
The most intuitive representation results from evaluating the nine scalars
\(\varLambda_{ij} = \mathbf{e}_i\cdot\bm{\varLambda}\mathbf{e}_j\) on a basis
\(\{\mathbf{e}_i\}\) spanning \(\mathbb{R}^3\) and forming a matrix, but
this is not recommended for implementation. It is well known that only
\emph{five} parameters are needed to uniquely and globally represent
elements of the rotation group \citep{stuelpnagel1964parametrization},
and in practice the non-unique double-cover\footnote{The unit quaternions are known as a ``double-cover'' because any single rotation is represented by two unit quaternions.} 
offered by the four \emph{Euler parameters} identifying a unit quaternion is generally the ideal choice.

In practical structural applications the impact of
the number of rotation parameters on computer memory requirements is often
inconsequential.
However, the size of the linear system that results from the tangent stiffness is important. 
This size is governed by the dimension of the \emph{variation} of a rotation, not the rotation itself. 
Because the rotation group \(\mathrm{SO}(3)\) is fundamentally a three-dimensional manifold, variations use only three parameters, regardless of the number of parameters used for the rotation variables.
Another consideration is the computational complexity of basic rotation operations for a given parameterization. 
The Euler parameters tend to be ideal in this regard, as quaternion operations generally incur the fewest trigonometric function evaluations.

\hypertarget{algebra}{%
\subsection{Algebra}\label{algebra}}

There is no useful way to combine elements of
\(\mathrm{SO}(3)\) associatively %which maintains
while maintaining the constraints of Equation~(\ref{eq:def-SO3}).
This means that when treating problems in terms of variables 
in \(\mathrm{SO}(3)\), addition is undefined, and one is left without the structure of a \emph{vector space}.
Fortunately, the collection \(\mathrm{SO}(3)\) as
defined preserves the weaker algebraic nature of a \emph{group}.
Although weaker than a vector space, this provides a very useful and
well-understood apparatus for algebraic manipulation (see e.g.
\cite{artin1991algebra}).

From this definition, it can also be shown that the group
\(\mathrm{SO}(3)\) forms a smooth manifold, allowing one to recover
constructs from calculus and work within the stronger structure of a
\emph{Lie group}.
The characterization as a Lie group is significant
because it allows operations in any tangent space, say
\(T_{\bm{R}}\mathrm{SO}(3)\) at an element \(\bm{R}\),
to be systematically related to the single tangent space
\(T_{\bm{1}}\mathrm{SO}(3)\) at the identity element \(\bm{1}\),
essentially representing the entire \emph{tangent bundle} \(\left\{T_{\scriptscriptstyle{\Lambda}}\,\mathrm{SO}(3)\mid \bm{\Lambda}\in\mathrm{SO}(3)\right\}\) %\FCF{???}
with a single vector space.
For this reason, the space \(T_{\bm{1}}\mathrm{SO}(3)\)
is given the special designation \(\mathfrak{so}(3)\) and is referred to
as the \emph{Lie algebra} of the group \(\mathrm{SO}(3)\).
Note that for problems in a vector space \(\mathcal{V}\),
this ordeal of tangent spaces %business \FCF{???!!!} 
is generally circumvented altogether by implicitly
identifying the tangent space \(T_x\mathcal{V}\) at each point \(x\) in
\(\mathcal{V}\) with the vector space \(\mathcal{V}\) itself.

The constraints in Equation (\ref{eq:def-SO3}) require that 
elements \(\bm{\Omega}\) of the Lie algebra \(\mathfrak{so}(3)\) behave like skew-symmetric linear operators. This immediately furnishes two convenient representations for these objects using either (1) vectors in \(\mathbb{R}^3\), or (2) the skew symmetric matrices in \(\mathbb{R}^{3 \times 3}\). 
Any vector \(\bm{\omega}\) in \(\mathbb{R}^3\) is associated
with a skew-symmetric matrix \(\bm{\Omega}\) using the \emph{hat} isomorphism
\((~\cdot~)^{\times}\). %\colon \mathbb{R}^3 \rightarrow \mathfrak{so}(3)\).
This is defined by the relation
\(\bm{\Omega}\bm{a} = \bm{\omega}^{\times} \bm{a}=\bm{\omega}\times \bm{a}\)
for any vector \(\bm{a}\) in \(\mathbb{R}^3\) and expressed in
coordinates as
\begin{equation}
\bm{\omega}^{\times}=\begin{pmatrix}
\omega_1 \\
\omega_2 \\
\omega_3
\end{pmatrix}^{\times}=\begin{pmatrix}
0 & -\omega_3 & \omega_2 \\
\omega_3 & 0 & -\omega_1 \\
-\omega_2 & \omega_1 & 0
\end{pmatrix} = \bm{\Omega}.
\label{eq:hat}\end{equation}
% Similarly, skew-symmetric operators are identified with a vector
% in \(\mathbb{R}^3\) using the \emph{vee} isomorphism
% \((~\cdot~)^{\vee}\). %\colon \mathfrak{so}(3)\rightarrow\mathbb{R}^3\). 
Similarly, given a skew-symmetric matrix \(\bm{\Omega}=\bm{\omega}^{\times}\),
the \emph{vee} isomorphism \((~\cdot~)^{\vee}\) is defined such that
\(\bm{\Omega}^{\vee}=\bm{\omega}\), which is expressed in
coordinates as: \begin{equation}
\bm{\Omega}^{\vee}= \begin{pmatrix}
0 & -\omega_3 & \omega_2 \\
\omega_3 & 0 & -\omega_1 \\
-\omega_2 & \omega_1 & 0
\end{pmatrix}^{\vee} = \begin{pmatrix}
\omega_1 \\
\omega_2 \\
\omega_3
\end{pmatrix}
= \bm{\omega}.
\label{eq:vee}\end{equation}

The inner product
\(\langle\cdot, \cdot\rangle\colon T_{\scriptscriptstyle{\Lambda}} \mathrm{SO}(3) \times T_{\scriptscriptstyle{\Lambda}} \mathrm{SO}(3) \rightarrow \mathbb{R}\)
is defined by (\cite{abraham1994foundations})

\begin{equation}
\left\langle\bm{\Omega}_{\bm{\Lambda}}, \bm{W}_{\bm{\Lambda}}\right\rangle
\triangleq \frac{1}{2} \operatorname{tr}\left[\bm{\Omega}_{\scriptscriptstyle{\Lambda}}^{\mathrm{t}} \bm{W}_{\scriptscriptstyle{\Lambda}}\right].
\label{eq:def-SO3-inner}\end{equation}
